\documentclass[11pt,a4paper]{article}

\usepackage[T1]{fontenc}
\usepackage{amsmath,amssymb,amsthm}
\usepackage[margin=2.6cm]{geometry}
\usepackage{mathtools}
\usepackage{array}
\usepackage{booktabs}

%% ---------------------------------------------------------------------------
%%  Bourbaki-style numbering: one running counter per section, shared by every
%%  numbered statement, with small-capitals headers.
%% ---------------------------------------------------------------------------
\newtheoremstyle{bourbaki}%
  {\topsep}{\topsep}{\itshape}{}{\scshape}{.}{ }{}
\newtheoremstyle{bourbakidef}%
  {\topsep}{\topsep}{\upshape}{}{\scshape}{.}{ }{}

\theoremstyle{bourbaki}
\newtheorem{thm}{Theorem}[section]
\newtheorem{prop}[thm]{Proposition}
\newtheorem{lem}[thm]{Lemma}
\newtheorem{cor}[thm]{Corollary}

\theoremstyle{bourbakidef}
\newtheorem{defn}[thm]{Definition}
\newtheorem{rem}[thm]{Remark}
\newtheorem{nota}[thm]{Notation}

\renewcommand{\qedsymbol}{$\blacksquare$}

\newcommand{\Z}{\mathbb{Z}}
\newcommand{\N}{\mathbb{N}}
\newcommand{\Q}{\mathbb{Q}}
\newcommand{\R}{\mathbb{R}}
\newcommand{\C}{\mathbb{C}}
\newcommand{\F}{\mathbb{F}}
\newcommand{\Lam}{\Lambda}
\newcommand{\eu}{e_{u}}
\newcommand{\ev}{e_{v}}
\DeclareMathOperator{\Mahler}{M}
\DeclareMathOperator{\Gal}{Gal}
\newcommand{\zmod}[1]{(\Z/#1\Z)^{\!\times}}

\title{A Fox--derivative cocycle for the Syracuse cycle problem:\\
       construction and the polynomials of short words}
\author{Leo Brunswic}
\date{21 June 2026}

\begin{document}
\maketitle

\begin{abstract}
We isolate the algebraic object underlying the cocycle calculus of the $3x+1$ problem,
stripped of the integer specialisation $(X,Y)\mapsto(2,3)$. On the free group $F_2=\langle
u,v\rangle$ with universal abelianisation character $\chi$, the composite $S=\chi\circ
\partial_u$ of the Fox derivative with $\chi$ is a $1$-cocycle valued in the Laurent ring
$\Lam=\Z[X^{\pm1},Y^{\pm1}]$. Section~1 develops $S$ from scratch: its closed form on
positive words, the cocycle law, the fact that its class in $\Lam/(1-\chi(w))$ is a
conjugacy invariant up to a monomial unit, and a unimodular reparametrisation that places
every word of type $(E,o)$ in the ring $R_d=\Z[t,t^{-1},s]/(s^d-1)$, $d=\gcd(E,o)$.
Section~2 specialises to \emph{coprime} type $\gcd(E,o)=1$, where $d=1$ and $R_1=\Z[t,t^{-1}]$
is a one-variable Laurent ring: we catalogue \emph{all} $1233$ cyclic classes with $E+o\le20$
($64$ types), computing for each the polynomial $P_w(t)$, its factorisation, Mahler measure
and (where small) Galois group. Each $P_w$ is a $0/1$ polynomial divisible by
$1+t+\cdots+t^{o-1}$ with a $\{-1,0,1\}$ (Littlewood) quotient. Two facts stand out. Every
coprime word with $E+o\le9$ is
cyclotomic, the first exception being type $(7,3)$ whose Mahler measure is the plastic number
$1.32472\ldots$ (smallest Pisot number); and the smallest Mahler measure $>1$ in the whole
range is \emph{Lehmer's number} $1.17628\ldots$, realised by an explicit word of type
$(12,7)$ whose cocycle factor is Lehmer's degree-$10$ polynomial. The four known integer cycles
appear as short words (values $-1,1,-5,-17$); the $-17$ cycle gives $\Phi_7$ times an
irreducible degree-$11$ factor of Mahler measure $1.427$. Section~3 places these in cohomology:
$R_w=\gcd(\chi\partial_u w,\chi\partial_v w)$ is the Alexander polynomial of the one-relator
group $\langle u,v\mid w\rangle$, the order of $H^2\!\big(G_w;\Lam/(1-\chi(w))\big)$, so the
Mahler/Lehmer phenomena are Alexander-polynomial phenomena.
\end{abstract}

\section{Construction}

\subsection*{The character and the Fox derivative}

\begin{nota}
$F_2=\langle u,v\rangle$ is the free group of rank $2$; $\Z[F_2]$ its integral group ring.
We write $\Lam=\Z[X^{\pm1},Y^{\pm1}]$ for the group ring of the abelianisation
$F_2^{\mathrm{ab}}=\Z^2$, identified with the Laurent polynomials in two variables.
\end{nota}

\begin{defn}\label{def:chi}
The \emph{universal character} is the ring homomorphism $\chi\colon \Z[F_2]\to\Lam$ extending
the group morphism $\chi(u)=X$, $\chi(v)=Y$. On a word $w$ it returns the abelianised
monomial $\chi(w)=X^{\eu(w)}Y^{\ev(w)}$, where $\eu,\ev$ are the (signed) exponent sums.
\end{defn}

\begin{defn}\label{def:fox}
The \emph{Fox derivative} $\partial_u\colon \Z[F_2]\to\Z[F_2]$ is the unique $\Z$-linear map
satisfying $\partial_u(1)=0$, $\partial_u(u)=1$, $\partial_u(v)=0$ and the Leibniz rule
\[
   \partial_u(ab)=\partial_u(a)+a\,\partial_u(b)\qquad(a,b\in F_2).
\]
From $\partial_u(g g^{-1})=0$ one gets $\partial_u(g^{-1})=-g^{-1}\partial_u(g)$.
\end{defn}

\begin{defn}\label{def:S}
The \emph{word cocycle} is the composite
\[
   S\;:=\;\chi\circ\partial_u\;\colon\;\Z[F_2]\longrightarrow\Lam .
\]
\end{defn}

\begin{prop}[Closed form on positive words]\label{prop:closed}
Let $w=x_1x_2\cdots x_E$ with each $x_k\in\{u,v\}$. Then
\[
   S(w)\;=\;\sum_{k:\,x_k=u} X^{\alpha_k}\,Y^{\beta_k},
   \qquad
   \alpha_k=\#\{i<k:\,x_i=u\},\quad \beta_k=\#\{i<k:\,x_i=v\}.
\]
In particular $S(w)$ is a sum of $\eu(w)$ monomials, each of total degree $\le E-1$, with
all coefficients equal to $1$.
\end{prop}

\begin{proof}
Iterating the Leibniz rule, $\partial_u(x_1\cdots x_E)=\sum_{k}x_1\cdots
x_{k-1}\,\partial_u(x_k)=\sum_{k:\,x_k=u}x_1\cdots x_{k-1}$, since $\partial_u(u)=1$ and
$\partial_u(v)=0$. Apply $\chi$ termwise: $\chi(x_1\cdots x_{k-1})=X^{\alpha_k}Y^{\beta_k}$.
\end{proof}

\begin{prop}[Cocycle law]\label{prop:cocycle}
$S(u)=1$, $S(v)=0$, and for all $a,b\in F_2$
\[
   S(ab)\;=\;S(a)+\chi(a)\,S(b).
\]
Equivalently $S\in Z^1\!\big(F_2,\Lam_\chi\big)$ is a $1$-cocycle for the action of $F_2$ on
$\Lam$ through $\chi$; these two properties determine $S$ uniquely.
\end{prop}

\begin{proof}
Apply $\chi$ to $\partial_u(ab)=\partial_u a+a\,\partial_u b$. Uniqueness: any map with
$S(u)=1,\,S(v)=0$ obeying the law is fixed on all words by induction on length.
\end{proof}

\subsection*{Type, closure, and conjugacy invariance}

\begin{defn}\label{def:type}
The \emph{type} of a positive word $w$ is the pair $(E,o)$ with $E=\eu(w)+\ev(w)=|w|$ (the
length) and $o=\eu(w)$ (the number of $u$'s). Thus $\chi(w)=X^{o}Y^{E-o}$ and
$S(w)$ has exactly $o$ terms.
\end{defn}

The interesting words are \emph{closed loops}: those for which the abelianised displacement
$\chi(w)$ is trivial, $\chi(w)=1$. We therefore reduce modulo the \emph{closure ideal}
$(1-\chi(w))\subset\Lam$. The following proposition is the reason this is the right
normalisation: it makes $S$ an invariant of the cyclic word.

\begin{prop}[Conjugacy invariance up to a unit]\label{prop:conj}
For $g\in F_2$,
\[
   S\!\big(g\,w\,g^{-1}\big)\;=\;\chi(g)\,S(w)\;+\;S(g)\,\big(1-\chi(w)\big).
\]
Consequently the class of $S(w)$ in the quotient $\Lam/(1-\chi(w))$ depends only on the
conjugacy class of $w$, up to multiplication by the monomial unit $\chi(g)$. Coboundaries
$w\mapsto c\,(1-\chi(w))$, $c\in\Lam$, act trivially on this class.
\end{prop}

\begin{proof}
Using Proposition~\ref{prop:cocycle} and $S(g^{-1})=-\chi(g)^{-1}S(g)$,
\[
   S(gwg^{-1})=S(g)+\chi(g)\big[S(w)+\chi(w)S(g^{-1})\big]
              =\chi(g)S(w)+S(g)\big(1-\chi(w)\big).
\]
Since $\chi$ is abelian, conjugate words share the same $\chi(w)$, hence the same ideal
$(1-\chi(w))$; modulo it the displayed identity reads $[S(gwg^{-1})]=\chi(g)\,[S(w)]$.
\end{proof}

Cyclic rotation of $w$ is conjugation by a prefix, so Proposition~\ref{prop:conj} says: the
rotations $S^{(j)}(w)$ all agree in $\Lam/(1-\chi(w))$ \emph{up to a monomial}. The natural
invariants of the cyclic word are therefore those insensitive to a monomial factor: the
factorisation type, the Newton polygon, the roots on the torus $|X|=|Y|=1$, the Mahler
measure, and the Galois group of the splitting field.

\subsection*{Reparametrisation: the ring $R_d$}

\begin{prop}[Normal form]\label{prop:reparam}
Let $w$ have type $(E,o)$, set $d=\gcd(E,o)=\gcd(o,E-o)$ and write $(o,E-o)=d\,(o',m')$ with
$(o',m')$ primitive. Choose $a,b\in\Z$ with $o'b-m'a=1$ and put
\[
   s=X^{o'}Y^{m'},\qquad t=X^{a}Y^{b}.
\]
Then $(s,t)$ is a $\Z$-basis of the monomial lattice, with inverse
\[
   X=s^{\,b}\,t^{-m'},\qquad Y=s^{-a}\,t^{\,o'},
\]
one has $\chi(w)=s^{d}$, and the reduction map induces a ring isomorphism
\[
   \Lam/\big(1-\chi(w)\big)\;\xrightarrow{\ \sim\ }\;R_d:=\Z[t,t^{-1},s]/(s^{d}-1)
   \;\cong\;\bigoplus_{j=0}^{d-1}\Z[t,t^{-1}]\,s^{\,j}.
\]
Under it $S(w)$ becomes a Laurent polynomial in $t$ whose coefficients are $\Z$-combinations
of $1,s,\dots,s^{d-1}$.
\end{prop}

\begin{proof}
The matrix $\left(\begin{smallmatrix}o'&m'\\ a&b\end{smallmatrix}\right)$ has determinant
$o'b-m'a=1$, hence lies in $\mathrm{GL}_2(\Z)$ and gives a monomial change of basis with the
stated inverse. Then $s^{d}=X^{o'd}Y^{m'd}=X^{o}Y^{E-o}=\chi(w)$, so the ideal $(1-\chi(w))$
becomes $(1-s^{d})=(s^{d}-1)$, and $t$ survives as a free Laurent variable.
\end{proof}

We normalise the (otherwise $(o',m')$-ambiguous) completion so that $t$ has smallest positive
$Y$-exponent, matching the shape $t=Y\,X^{*}Y^{*}$; for $m'=0$ this gives $t=Y$.

\begin{rem}[The study programme]\label{rem:program}
$R_d$ is free of rank $d$ over $\Z[t^{\pm1}]$, with $s$ a formal $d$-th root of unity.
Specialising $s\mapsto\zeta$ over the $d$-th roots of unity sends $S(w)$ to a Laurent
polynomial $S(w)\big|_{s=\zeta}\in\Z[\zeta][t^{\pm1}]$. The programme of this approach is to
study, systematically and across words, the \emph{roots, factorisations, Galois groups and
Mahler measures} of these polynomials. Specialising further to $(X,Y)\mapsto(2,3)$ recovers
the integer cocycle of the \texttt{cocycle-combine} program, where $S(w)/(\chi(w)-1)$ is the
rational fixed point (cycle start) of the word; cyclotomic factors and small Mahler measure
are exactly what force the arithmetic ceilings there. Here we keep $(X,Y)$ formal and read
off the pure algebra.
\end{rem}


\section{The coprime catalogue: $\gcd(E,o)=1$, $E+o\le 20$}

When $\gcd(E,o)=1$ we have $d=1$, so $R_1=\Z[t,t^{-1}]$ and \emph{every cyclic word is a
single Laurent polynomial in one variable}. Concretely the closure $\chi(w)=1$ is imposed by
the substitution of Proposition~\ref{prop:reparam} with $s=1$,
\[
   X\longmapsto t^{-(E-o)},\qquad Y\longmapsto t^{\,o},
\]
under which a monomial $X^{\alpha}Y^{\beta}$ of $S(w)$ becomes $t^{\,\beta o-\alpha(E-o)}$.
Clearing the lowest power of $t$ gives an integer polynomial $P_w(t)\in\Z[t]$, well defined up
to $t\mapsto t^{-1}$ and a power of $t$ (Proposition~\ref{prop:conj}); e.g.\ $uuv\mapsto t+1$
and $uuuv\mapsto t^2+t+1=\Phi_3$. The relevant invariants are those of $P_w$: its factorisation
into cyclotomic and non-cyclotomic parts, its Mahler measure, and the Galois groups of its
irreducible factors. (For general $d$ one keeps $s$: e.g.\ $uvuv\mapsto 1+s$ in $R_2$ and
$uuu\mapsto1+s+s^2$ in $R_3$; all $10$ classes with $E+o\le6$, any $d$, are cyclotomic.)

We enumerate \emph{every} coprime cyclic class with $E+o\le20$: $1233$ classes across $64$
types. All data are produced by \texttt{gen\_table.py} (the cocycle law of
Proposition~\ref{prop:cocycle} is re-verified on every split of every word up to length~$6$)
and recorded in full in \texttt{coprime\_catalogue.csv} (which lists each $P_w$ and its
Littlewood core $R_w$ of Proposition~\ref{prop:phio} explicitly); Table~\ref{tab:summary}
aggregates by type and Table~\ref{tab:highlights} lists the arithmetically extreme classes.

\subsection*{The four known cycles}

The four known cycles of the $3n+1$ map on $\Z$ are short words of the accelerated $T$-map
($u:x\mapsto(3x+1)/2$, $v:x\mapsto x/2$): the parity sequence of a periodic orbit is a positive
word $w$, its fixed point $F_w(x)=x$ is the cycle value, and the type $(E,o)$ gives
$D=2^{E}-3^{o}$. They are $\mathtt{u}$ ($\{-1\}$), $\mathtt{uv}$ ($\{1,2\}$), $\mathtt{uuv}$
(the $-5$ cycle), and $\mathtt{uuuuvuuuvvv}$ (the $-17$ cycle, $7$ odd steps); all four are
coprime. The Mahler measures and $D=2^E-3^o$ are exact; the fixed points reproduce the cycle
values (checked in \texttt{gen\_table.py}).

\begin{prop}[The cyclic factor and the Littlewood core]\label{prop:phio}
Let $w$ be a cyclic word of coprime type. Then $P_w$ is a $0/1$ polynomial with exactly $o$
terms, the full cyclotomic factor
\[
   g_o(t)\;:=\;\frac{t^{o}-1}{t-1}\;=\;1+t+\cdots+t^{o-1}\;=\!\!\prod_{d\mid o,\ d>1}\!\!\Phi_d(t)
\]
divides $P_w(t)$, and the quotient $R_w:=P_w/g_o$ has all coefficients in $\{-1,0,1\}$.
\end{prop}

\begin{proof}
By Proposition~\ref{prop:closed} (the $k$-th $u$ has $\alpha_k=k$ earlier $u$'s) the exponents
$e_k=o\beta_k-(E-o)k$ satisfy $e_k\equiv-(E-o)k\pmod o$; since $\gcd(E-o,o)=\gcd(E,o)=1$ the map
$k\mapsto-(E-o)k$ is a bijection of $\Z/o\Z$, so $\{e_k\bmod o\}$ is a complete residue system.
In particular the $e_k$ are pairwise distinct, so $P_w=\sum_k t^{e_k}$ is $0/1$ with $o$ terms;
and reducing each $t^{e_k}$ modulo $t^{o}-1$ gives $P_w\equiv\sum_{j=0}^{o-1}t^{j}=g_o\pmod{t^{o}-1}$,
so $P_w$ vanishes at every $o$-th root of unity except $1$, i.e.\ $g_o\mid P_w$.

For the quotient, write $[P]_m$ for the coefficient of $t^m$ and let $e(\rho)$ be the unique
exponent of $P_w$ in the class $\rho\bmod o$. From $R_w\,(t^{o}-1)=P_w\,(t-1)$ one gets
$R_n=R_{n+o}+[P_w(t-1)]_{n+o}$; summing downward from $R_{+\infty}=0$,
\[
   R_n=\sum_{i\ge1}\bigl([P_w]_{n+oi-1}-[P_w]_{n+oi}\bigr)
      =\mathbf 1\!\left[e\bigl((n{-}1)\bmod o\bigr)>n{-}1\right]
      -\mathbf 1\!\left[e\bigl(n\bmod o\bigr)>n\right],
\]
a difference of two $0/1$ indicators (one exponent per residue class), hence in $\{-1,0,1\}$.
\end{proof}

\begin{cor}[Normalisation]\label{cor:Rone}
$R_w(1)=1$; equivalently $R_w=1+(t-1)\,Q_w$, where
$Q_w(t)=\sum_{n\ge0}\mathbf 1\!\left[e(n\bmod o)>n\right]t^{n}$ is itself a $0/1$ polynomial. Hence
\[
   P_w \;=\; g_o\,\bigl(1+(t-1)\,Q_w\bigr),\qquad g_o,\,Q_w\in\{0,1\}[t],\quad R_w\in\{-1,0,1\}[t].
\]
\end{cor}

\begin{proof}
$P_w$ has exactly $o$ terms, so $R_w(1)=P_w(1)/g_o(1)=o/o=1$. With
$f(n):=\mathbf 1[e(n\bmod o)>n]$ the proof of Proposition~\ref{prop:phio} gives $R_n=f(n-1)-f(n)$,
whose tail sums telescope: $\bigl[(R_w-1)/(t-1)\bigr]_n=\sum_{i>n}R_i=f(n)\in\{0,1\}$.
\end{proof}

\begin{table}[h]\centering\small\renewcommand{\arraystretch}{1.25}
\begin{tabular}{@{}llccc@{}}
\toprule
word $w$ & cycle & $(E,o)$ & $D=2^{E}-3^{o}$ & $\mathrm{M}(P_w)$\\
\midrule
$\mathtt{u}$ & $\{-1\}$ & $(1,1)$ & $-1$ & $1$\\
$\mathtt{uv}$ & $\{1\}$ & $(2,1)$ & $1$ & $1$\\
$\mathtt{uuv}$ & $\{-5\}$ & $(3,2)$ & $-1$ & $1$\\
$\mathtt{uuuuvuuuvvv}$ & $\{-17\}$ & $(11,7)$ & $-139$ & $1.42747$\\
\bottomrule
\end{tabular}
\caption{The four known cycles of the $3n+1$ map on $\mathbb{Z}$, read as words of the accelerated $T$-map ($u:x\mapsto(3x+1)/2$, $v:x\mapsto x/2$): the cycle value $x$ (fixed point of the word map), the type, $D=2^E-3^o$, and the Mahler measure of $P_w$.}\label{tab:cycles}
\end{table}
\begin{align*}
P_{-5}(t) &= t + 1\\
 &= \Phi_{2}.\\[4pt]
P_{-17}(t) &= t^{17} + t^{13} + t^{9} + t^{8} + t^{5} + t^{4} + 1\\
 &= \Phi_{7}\,\bigl(t^{11} - t^{10} + t^{7} - t^{6} + t^{4} - t + 1\bigr).
\end{align*}


The two trivial cycles give $P_w=1$; the $-5$ cycle gives $P_w=\Phi_2$ (purely cyclotomic);
and the $-17$ cycle is the first to carry arithmetic, $P_w=g_7\,R_w$ with $g_7=\Phi_7$ (here
$o=7$ is prime) and Littlewood core $R_w=t^{11}-t^{10}+t^{7}-t^{6}+t^{4}-t+1$, irreducible and
non-reciprocal, of Mahler measure $1.42747\ldots$ (six roots off the unit circle). The
cyclotomic factor $g_o$ is forced by Proposition~\ref{prop:phio}; the core $R_w=P_w/g_o$ is the
genuine invariant attached to the cycle (Corollary~\ref{cor:Rone} writes it as $1+(t-1)Q_w$ with
$Q_w=t^{10}+t^{6}+t^{3}+t^{2}+t$).

\begin{prop}[Coprime cyclotomic frontier]\label{prop:cofrontier}
Every coprime cyclic word with $E+o\le 9$ has $\Mahler(P_w)=1$; i.e.\ $P_w$ is a product of
cyclotomic polynomials. The first non-cyclotomic class occurs at $E+o=10$, type $(7,3)$: for
$w=\mathtt{uuvuvvv}$,
\[
   P_w \;=\; \Phi_3(t)\,\bigl(t^{3}-t^{2}+1\bigr),\qquad
   \Mahler(P_w)=1.324718\ldots,
\]
the plastic number (the smallest Pisot number, real root of $x^{3}-x-1$); its non-cyclotomic
factor has Galois group $S_3$.
\end{prop}

\begin{prop}[Lehmer's number is realised]\label{prop:lehmer}
Across the $1233$ classes the Mahler measures lie in $[1,\,1.65378]$, with $274$ cyclotomic and
$959$ non-cyclotomic. The smallest measure exceeding $1$ is \emph{Lehmer's number}
$\Mahler=1.1762808\ldots$, attained (among others) by $w=\mathtt{uuuvuuvvvuuv}$ of type
$(12,7)$, whose non-cyclotomic factor is Lehmer's polynomial
\[
   t^{10}-t^{9}+t^{7}-t^{6}+t^{5}-t^{4}+t^{3}-t+1 .
\]
This is the smallest Mahler measure $>1$ known for any integer polynomial (Lehmer's problem);
that it is realised by a short Syracuse word is the catalogue's most striking feature.
\end{prop}

\begin{table}[h]
\centering\footnotesize
\renewcommand{\arraystretch}{1.1}
\setlength{\tabcolsep}{4.5pt}
\begin{tabular}{@{}cccccc@{\qquad}cccccc@{}}
\toprule
$(E,o)$ & cl. & cy. & nc. & min & max & $(E,o)$ & cl. & cy. & nc. & min & max\\
\midrule
$(1,1)$ & 1 & 1 & 0 & $--$ & $1.0000$ & $(11,1)$ & 1 & 1 & 0 & $--$ & $1.0000$ \\
$(2,1)$ & 1 & 1 & 0 & $--$ & $1.0000$ & $(11,2)$ & 5 & 5 & 0 & $--$ & $1.0000$ \\
$(3,1)$ & 1 & 1 & 0 & $--$ & $1.0000$ & $(11,3)$ & 15 & 5 & 10 & $1.3247$ & $1.3716$ \\
$(3,2)$ & 1 & 1 & 0 & $--$ & $1.0000$ & $(11,4)$ & 30 & 10 & 20 & $1.3247$ & $1.4864$ \\
$(4,1)$ & 1 & 1 & 0 & $--$ & $1.0000$ & $(11,5)$ & 42 & 8 & 34 & $1.2612$ & $1.5367$ \\
$(4,3)$ & 1 & 1 & 0 & $--$ & $1.0000$ & $(11,6)$ & 42 & 8 & 34 & $1.2612$ & $1.5367$ \\
$(5,1)$ & 1 & 1 & 0 & $--$ & $1.0000$ & $(11,7)$ & 30 & 10 & 20 & $1.3247$ & $1.4864$ \\
$(5,2)$ & 2 & 2 & 0 & $--$ & $1.0000$ & $(11,8)$ & 15 & 5 & 10 & $1.3247$ & $1.3716$ \\
$(5,3)$ & 2 & 2 & 0 & $--$ & $1.0000$ & $(11,9)$ & 5 & 5 & 0 & $--$ & $1.0000$ \\
$(5,4)$ & 1 & 1 & 0 & $--$ & $1.0000$ & $(12,1)$ & 1 & 1 & 0 & $--$ & $1.0000$ \\
$(6,1)$ & 1 & 1 & 0 & $--$ & $1.0000$ & $(12,5)$ & 66 & 5 & 61 & $1.1763$ & $1.5586$ \\
$(6,5)$ & 1 & 1 & 0 & $--$ & $1.0000$ & $(12,7)$ & 66 & 5 & 61 & $1.1763$ & $1.5586$ \\
$(7,1)$ & 1 & 1 & 0 & $--$ & $1.0000$ & $(13,1)$ & 1 & 1 & 0 & $--$ & $1.0000$ \\
$(7,2)$ & 3 & 3 & 0 & $--$ & $1.0000$ & $(13,2)$ & 6 & 6 & 0 & $--$ & $1.0000$ \\
$(7,3)$ & 5 & 3 & 2 & $1.3247$ & $1.3247$ & $(13,3)$ & 22 & 6 & 16 & $1.3247$ & $1.3756$ \\
$(7,4)$ & 5 & 3 & 2 & $1.3247$ & $1.3247$ & $(13,4)$ & 55 & 15 & 40 & $1.3247$ & $1.5189$ \\
$(7,5)$ & 3 & 3 & 0 & $--$ & $1.0000$ & $(13,5)$ & 99 & 10 & 89 & $1.1763$ & $1.6001$ \\
$(7,6)$ & 1 & 1 & 0 & $--$ & $1.0000$ & $(13,6)$ & 132 & 14 & 118 & $1.2353$ & $1.6538$ \\
$(8,1)$ & 1 & 1 & 0 & $--$ & $1.0000$ & $(13,7)$ & 132 & 14 & 118 & $1.2353$ & $1.6538$ \\
$(8,3)$ & 7 & 3 & 4 & $1.3247$ & $1.3497$ & $(14,1)$ & 1 & 1 & 0 & $--$ & $1.0000$ \\
$(8,5)$ & 7 & 3 & 4 & $1.3247$ & $1.3497$ & $(14,3)$ & 26 & 6 & 20 & $1.3247$ & $1.3768$ \\
$(8,7)$ & 1 & 1 & 0 & $--$ & $1.0000$ & $(14,5)$ & 143 & 9 & 134 & $1.2164$ & $1.6038$ \\
$(9,1)$ & 1 & 1 & 0 & $--$ & $1.0000$ & $(15,1)$ & 1 & 1 & 0 & $--$ & $1.0000$ \\
$(9,2)$ & 4 & 4 & 0 & $--$ & $1.0000$ & $(15,2)$ & 7 & 7 & 0 & $--$ & $1.0000$ \\
$(9,4)$ & 14 & 6 & 8 & $1.3247$ & $1.4458$ & $(15,4)$ & 91 & 21 & 70 & $1.3247$ & $1.5298$ \\
$(9,5)$ & 14 & 6 & 8 & $1.3247$ & $1.4458$ & $(16,1)$ & 1 & 1 & 0 & $--$ & $1.0000$ \\
$(9,7)$ & 4 & 4 & 0 & $--$ & $1.0000$ & $(16,3)$ & 35 & 7 & 28 & $1.3247$ & $1.3782$ \\
$(9,8)$ & 1 & 1 & 0 & $--$ & $1.0000$ & $(17,1)$ & 1 & 1 & 0 & $--$ & $1.0000$ \\
$(10,1)$ & 1 & 1 & 0 & $--$ & $1.0000$ & $(17,2)$ & 8 & 8 & 0 & $--$ & $1.0000$ \\
$(10,3)$ & 12 & 4 & 8 & $1.3247$ & $1.3679$ & $(17,3)$ & 40 & 8 & 32 & $1.3247$ & $1.3787$ \\
$(10,7)$ & 12 & 4 & 8 & $1.3247$ & $1.3679$ & $(18,1)$ & 1 & 1 & 0 & $--$ & $1.0000$ \\
$(10,9)$ & 1 & 1 & 0 & $--$ & $1.0000$ & $(19,1)$ & 1 & 1 & 0 & $--$ & $1.0000$ \\
\bottomrule
\end{tabular}
\caption{All $64$ coprime types ($\gcd(E,o)=1$, $E+o\le20$), in two side-by-side blocks
of $32$: number of cyclic classes (cl.), how many are cyclotomic (cy.) and non-cyclotomic
(nc.), the smallest Mahler measure $>1$ in the type (min; ``--'' if none), and the largest
Mahler measure (max).}\label{tab:summary}
\end{table}


\begin{table}[h]
\centering\small
\renewcommand{\arraystretch}{1.2}
\begin{tabular}{@{}llcccl@{}}
\toprule
$w$ & $(E,o)$ & $\deg P_w$ & $\deg_{\mathrm{nc}}$ & $\mathrm{M}$ & class\\
\midrule
$\mathtt{uuuvuuvvvuuv}$ & $(12,7)$ & $16$ & $10$ & $1.17628$ & Lehmer \\
$\mathtt{uuuvvuvvvuvv}$ & $(12,5)$ & $14$ & $10$ & $1.17628$ & Lehmer \\
$\mathtt{uuvuvvvuvvvuv}$ & $(13,5)$ & $14$ & $10$ & $1.17628$ & Lehmer \\
$\mathtt{uuuvvuvvvvvuvv}$ & $(14,5)$ & $18$ & $10$ & $1.21639$ & Salem \\
$\mathtt{uuuuvvuvvvuvv}$ & $(13,6)$ & $21$ & $16$ & $1.23526$ & 2 ext. \\
$\mathtt{uuuvuuvvvvuuv}$ & $(13,7)$ & $22$ & $16$ & $1.23526$ & 2 ext. \\
$\mathtt{uuvuuvvvuvvv}$ & $(12,5)$ & $16$ & $12$ & $1.25105$ & 2 ext. \\
$\mathtt{uuuvvvuvvuvvv}$ & $(13,5)$ & $16$ & $12$ & $1.25105$ & 2 ext. \\
$\mathtt{uuvvuvvuuvvvvv}$ & $(14,5)$ & $16$ & $12$ & $1.25105$ & 2 ext. \\
$\mathtt{uuuvuuuvvuvv}$ & $(12,7)$ & $18$ & $12$ & $1.25105$ & 2 ext. \\
$\mathtt{uuuuvuvvvuv}$ & $(11,6)$ & $15$ & $10$ & $1.26123$ & Salem \\
$\mathtt{uuuvuvvvvuv}$ & $(11,5)$ & $14$ & $10$ & $1.26123$ & Salem \\
$\mathtt{uuvuuvvvvuvvvv}$ & $(14,5)$ & $22$ & $18$ & $1.26633$ & 3 ext. \\
$\mathtt{uuuvuvuuuvvv}$ & $(12,7)$ & $16$ & $10$ & $1.28358$ & 2 ext. \\
$\mathtt{uuvuuvuuvvuvv}$ & $(13,7)$ & $16$ & $10$ & $1.28358$ & 2 ext. \\
$\mathtt{uuuuvvuvuvv}$ & $(11,6)$ & $15$ & $10$ & $1.28358$ & 2 ext. \\
$\mathtt{uuvuuvvuvvuvv}$ & $(13,6)$ & $15$ & $10$ & $1.28358$ & 2 ext. \\
$\mathtt{uuvuvuuvvvv}$ & $(11,5)$ & $14$ & $10$ & $1.28358$ & 2 ext. \\
$\mathtt{uuuvvvuvuvvv}$ & $(12,5)$ & $14$ & $10$ & $1.28358$ & 2 ext. \\
$\mathtt{uuvvuuvvvuvvv}$ & $(13,5)$ & $14$ & $10$ & $1.28358$ & 2 ext. \\
$\mathtt{uuvuvuuvvvvvvv}$ & $(14,5)$ & $26$ & $22$ & $1.28915$ & 5 ext. \\
$\mathtt{uuuuuvuvvvuv}$ & $(12,7)$ & $20$ & $12$ & $1.30443$ & 2 ext. \\
$\mathtt{uuuvuvvvvvuv}$ & $(12,5)$ & $18$ & $12$ & $1.30443$ & 2 ext. \\
$\mathtt{uuuvuvvvvvvvuv}$ & $(14,5)$ & $26$ & $20$ & $1.30499$ & 4 ext. \\
\bottomrule
\end{tabular}
\caption{The $24$ coprime classes of smallest Mahler measure $>1$: word, type, degree of
$P_w$, degree of its non-cyclotomic factor, Mahler measure, and the nature of that factor
(Lehmer's polynomial; a Salem number $=$ reciprocal with one root off the unit circle; a
Pisot number; or ``$k$ ext.'' $=$ $k$ roots outside the circle).}\label{tab:highlights}
\end{table}


\begin{rem}
The cyclotomic part of $P_w$ collects roots of unity; the non-cyclotomic part carries the
arithmetic. Small Mahler measure means a Salem or Pisot factor, and these recur across many
words and types: within $E+o\le20$ Lehmer's value $1.17628$ is realised by several words, and
the next measures attained are $1.21639,\,1.23526,\,1.25105,\ldots$ Specialising $t$ further to
$(X,Y)=(2,3)$ feeds exactly these factors into
the divisibility ceilings of \texttt{cocycle-combine}: a Salem/Lehmer factor is a
\emph{near-cyclotomic} obstruction, with Mahler measure just above $1$. Extending the catalogue
past $E+o=20$, and pinning down which words realise the successive small Salem numbers, is the
natural continuation.
\end{rem}

\subsection*{The map $w\mapsto R_w$}

The core is a true invariant of the cyclic word, with two reflection symmetries. Write $w^{R}$
for the reversed word and $\bar w$ for the complement (swap $u\leftrightarrow v$, which sends
type $(E,o)$ to $(E,E-o)$).

\begin{prop}[Reflections and fibres]\label{prop:cores}
$R_{w^{R}}$ and $R_{\bar w}$ both equal the reciprocal $t^{\deg R_w}R_w(1/t)$. Hence
$R_{\overline{w^{R}}}=R_w$: the involution $w\mapsto\overline{w^{R}}$, of type
$(E,o)\mapsto(E,E-o)$, preserves the core, so the cores occur in fibres closed under it. Moreover
$R_w(0)=R_w(1)=1$.
\end{prop}

\begin{proof}
$R_w(0)=P_w(0)/g_o(0)=1$ (both have constant term $1$) and $R_w(1)=1$ by
Corollary~\ref{cor:Rone}. Reversal exchanges prefixes with suffixes and complement exchanges
$(o,E-o)$; on the reduced cocycle each acts by $t\mapsto t^{-1}$ --- for the complement via the
Fox fundamental formula $(X-1)S_u(w)+(Y-1)S_v(w)=\chi(w)-1$ together with $S_u(\bar w)=S_v(w)$
--- and $g_o$ is reciprocal. Verified on all coprime necklaces with $E\le16$.
\end{proof}

The map is very far from injective: the $4506$ coprime necklaces with $E\le16$ realise only
$1529$ distinct cores ($818$ up to reciprocal). The simple cores have enormous fibres ---
$R_w=1$ for $80$ necklaces, $R_w=\Phi_6=t^2-t+1$ for $50$ --- and the $55$ \emph{cyclotomic}
cores (Mahler $1$) have mean fibre $7.7$ against $2.8$ for the $1474$ non-cyclotomic ones. So the
arithmetic of a word concentrates in a rare non-cyclotomic core shared by only a few necklaces.
Table~\ref{tab:cores} lists the small-degree cores; the full distinct-core catalogue is
\texttt{R\_cores.csv}.

\begin{table}[h]\centering\small\renewcommand{\arraystretch}{1.2}
\begin{tabular}{@{}lcccc@{}}
\toprule
core $R_w(t)$ & $\deg$ & $\mathrm{M}$ & cyclotomic & fibre\\
\midrule
$1$ & $0$ & $1$ & yes & $80$\\
$t^{2} - t + 1$ & $2$ & $1$ & yes & $50$\\
$t^{3} - t^{2} + 1$ & $3$ & $1.3247$ & no & $38$\\
$t^{3} - t + 1$ & $3$ & $1.3247$ & no & $38$\\
$t^{4} - t^{3} + t^{2} - t + 1$ & $4$ & $1$ & yes & $34$\\
$t^{4} - t^{3} + 1$ & $4$ & $1.4013$ & no & $24$\\
$t^{4} - t^{2} + 1$ & $4$ & $1$ & yes & $24$\\
$t^{4} - t + 1$ & $4$ & $1.4013$ & no & $24$\\
\bottomrule
\end{tabular}
\caption{The distinct cores $R_w$ of degree $\le4$, over all $4506$ coprime necklaces with $E\le16$ ($1529$ distinct cores, $55$ of them cyclotomic). Each recurs across a whole fibre of necklaces; cyclotomic cores dominate. Cores come in reciprocal pairs (e.g.\ $t^3-t^2+1$ and $t^3-t+1$), self-reciprocal ones single.}\label{tab:cores}
\end{table}


\paragraph{How $R_w$ splits.} For coprime type the ambient ring $\Lam/(1-\chi(w))=\Z[t,t^{-1}]$
is a domain, so ``splitting'' is the factorisation of $R_w$ into cyclotomic factors times a
non-cyclotomic Salem/Pisot factor --- the latter absent exactly for the $55$ cyclotomic cores.
For \emph{non-coprime} type, $d=\gcd(E,o)>1$, the ring itself splits: by the Chinese remainder
theorem (over $\Q$)
\[
   \Q[t,t^{-1},s]/(s^{d}-1)\;\cong\;\prod_{e\mid d}\Q(\zeta_e)[t,t^{-1}],
\]
one factor per divisor $e$ of $d$, and the reduced cocycle distributes over the factors; the
$e=1$ component recovers the present $R_w$, while the $e>1$ components carry genuinely new
$\zeta_e$-twisted data. The non-coprime stratum is the natural next target.


\section{Cohomological frame: $R_w$ as an Alexander polynomial}

We have produced three objects --- the cocycle $S$, the reduced polynomial $P_w$, and the core
$R_w$. This section fixes their cohomological meaning: $S$ is a $1$-cocycle on $F_2$, $P_w$ is
its restriction to the loop $\langle w\rangle$, and $R_w$ is the \emph{Alexander polynomial} of
the one-relator group $\langle u,v\mid w\rangle$ --- the order of an $H^2$.

\subsection*{The cocycle and its restriction}

By Proposition~\ref{prop:cocycle}, $S=\chi\circ\partial_u$ is a $1$-cocycle, so
$[S]\in H^1(F_2;\Lam)$ for the coefficients $\Lam$ with $F_2$ acting through $\chi$. As $F_2$ is
free it has cohomological dimension $1$, whence $H^n(F_2;-)=0$ for $n\ge2$: there is no $H^2$
over $F_2$ itself, and the relevant class lives in $H^1$.

Restrict to the cyclic subgroup $\langle w\rangle\cong\Z$. For $\Z$ acting on a module $M$
through $\tau:=\chi(w)$ one has $H^1(\Z;M)=M/(\tau-1)M$ (the coinvariants), a $1$-cocycle being
determined by its value on the generator modulo $(\tau-1)M$. Hence:

\begin{prop}\label{prop:restrict}
The restriction of $[S]$ to $\langle w\rangle$ is
\[
   \operatorname{res}_{\langle w\rangle}[S]\;=\;[\,S(w)\,]\;=\;P_w
   \ \in\ H^1(\langle w\rangle;\Lam)=\Lam/(1-\chi(w)),
\]
the class of $S(w)$ modulo the closure ideal (Proposition~\ref{prop:conj}). So $P_w$ is the
cocycle class evaluated on the loop $w$.
\end{prop}

\subsection*{The Alexander matrix of the relator}

Alongside $S=\chi\circ\partial_u$ put the companion cocycle $T:=\chi\circ\partial_v$; on positive
words $T(w)=\sum_{k:\,x_k=v}X^{\alpha_k}Y^{\beta_k}$ (Proposition~\ref{prop:closed} with $u,v$
exchanged).

\begin{lem}[Crowell--Fox fundamental formula]\label{lem:cf}
$\partial_u w\,(u-1)+\partial_v w\,(v-1)=w-1$ in $\Z[F_2]$; applying $\chi$,
\[
   (X-1)\,S(w)+(Y-1)\,T(w)=\chi(w)-1 ,
\]
so $(X-1)S(w)+(Y-1)T(w)\equiv0$ in $\Lam/(1-\chi(w))$. (Certified in \texttt{gen\_table.py}.)
\end{lem}

\subsection*{The one-relator group and $H^2$}

Let $G_w:=\langle u,v\mid w\rangle$. For coprime type the necklace is aperiodic, so $w$ is not a
proper power in $F_2$; by Lyndon's theorem the presentation $2$-complex is aspherical, a
$K(G_w,1)$. The character $\chi$ descends to $G_w$ exactly in $M:=\Lam/(1-\chi(w))$, where
$\chi(w)=1$, making $M$ a $G_w$-module (for coprime type the image of $\chi$ is the infinite
cyclic group $\langle t\rangle$, so $M=\Z[t,t^{-1}]=\Z[\Z]$). The cellular cochain complex of the
presentation complex (one $0$-cell, two $1$-cells $u,v$, one $2$-cell $w$),
\[
   M\;\xrightarrow{\ d^0\ }\;M^2\;\xrightarrow{\ d^1\ }\;M,\qquad
   d^1(m_u,m_v)=S(w)\,m_u+T(w)\,m_v,
\]
has top map the transpose of the Fox Jacobian $\partial_2=(\partial_u w,\partial_v w)$ --- the
\emph{Alexander matrix} of the relator. Hence
\[
   H^2(G_w;M)=\operatorname{coker}(d^1)=M/\big(S(w),T(w)\big).
\]

\subsection*{$R_w$ is the Alexander polynomial}

\begin{thm}\label{thm:alex}
For a coprime word $w$, over $\Q[t^{\pm1}]$,
\[
   \gcd\big(S(w),\,T(w)\big)\;=\;R_w .
\]
Consequently the elementary ideal $(S(w),T(w))\otimes\Q=(R_w)$, and
\[
   \Q\otimes H^2(G_w;M)\;\cong\;\Q[t^{\pm1}]/(R_w).
\]
Thus $R_w$ is the \emph{Alexander polynomial} $\Delta_{(G_w,\chi)}$ --- the order of the
(rational) Alexander module --- and $\deg R_w=\dim_\Q\big(\Q\otimes H^2(G_w;M)\big)$.
\end{thm}

\begin{proof}
Write $P_w=S(w)$, $T_w=T(w)$ in $M=\Z[t^{\pm1}]$, and $g_k:=1+t+\cdots+t^{k-1}=(t^k-1)/(t-1)$.
Clearing denominators in Lemma~\ref{lem:cf} (with $X=t^{-(E-o)}$, $Y=t^{o}$) gives
$(t^{E-o}-1)P_w=t^{E-o}(t^{o}-1)T_w$, i.e.\ $g_{E-o}\,P_w=t^{E-o}g_o\,T_w$. By
Corollary~\ref{cor:Rone}, $P_w=g_o R_w$; substituting and cancelling $g_o$ yields
$T_w=t^{-(E-o)}g_{E-o}R_w$ --- a unit times $g_{E-o}R_w$. Therefore
\[
   \gcd(P_w,T_w)=R_w\cdot\gcd(g_o,g_{E-o}).
\]
The common roots of $g_o$ and $g_{E-o}$ would be common nontrivial $o$-th and $(E-o)$-th roots
of unity; since $\gcd(o,E-o)=\gcd(E,o)=1$ there are none, so $\gcd(g_o,g_{E-o})=1$ and
$\gcd(P_w,T_w)=R_w$. Over the principal ideal domain $\Q[t^{\pm1}]$ the ideal $(P_w,T_w)$ is then
$(R_w)$, whence $\Q\otimes H^2(G_w;M)=\Q[t^{\pm1}]/(R_w)$.
\end{proof}

\subsection*{Dictionary}

\begin{rem}\label{rem:dict}
Theorem~\ref{thm:alex} reads the whole catalogue as Alexander-module data.
\begin{itemize}
\item[(i)] $\Mahler(R_w)$ is the Mahler measure of an Alexander polynomial. Lehmer's number is
the Mahler measure of the Alexander polynomial of the $(-2,3,7)$-pretzel knot; its appearance at
type $(12,7)$ (Proposition~\ref{prop:lehmer}) is the cocycle realisation of that classical fact.
\item[(ii)] The reflection symmetry of Proposition~\ref{prop:cores} is \emph{external}:
complement-then-reverse exchanges $\partial_u\leftrightarrow\partial_v$ (i.e.\ $S\leftrightarrow T$)
and $t\mapsto t^{-1}$, and as $\gcd$ is symmetric it carries $R_w$ to the same core. Individual
cores need \emph{not} be self-reciprocal --- unlike a knot's $\Delta$ --- see
Remark~\ref{rem:duality}.
\item[(iii)] $P_w=g_o R_w$ (Corollary~\ref{cor:Rone}) is the splitting ``(forced cyclotomic
part, killed in passing from $H^1(\langle w\rangle)$ to $H^2(G_w)$) $\times$ (Alexander core
$R_w$)''.
\item[(iv)] The four cycles of Section~2 carry their Alexander polynomials: $R_w=1$ for
$\{-1\},\{1\}$; $R_w=\Phi_2$ for $-5$; and $R_w=t^{11}-t^{10}+t^{7}-t^{6}+t^{4}-t+1$ for $-17$
(so $\dim_\Q H^2=11$).
\end{itemize}
\end{rem}

\subsection*{The Fox--Lyndon complex, torsion, and the deeper strata}

\begin{rem}[The resolution and its torsion]\label{rem:torsion}
The objects above are the low-degree pieces of one complex. The presentation $2$-complex of
$G_w$ being aspherical, its universal cover gives a length-$2$ free resolution of $\Z$ over
$\Z[G_w]$,
\[
   0\to \Z[G_w]\xrightarrow{\,(\partial_u w,\,\partial_v w)\,}\Z[G_w]^2
        \xrightarrow{\,(u-1,\,v-1)\,}\Z[G_w]\xrightarrow{\ \varepsilon\ }\Z\to0,
\]
with the Fox Jacobian as its only syzygy (Fox's free calculus; Lyndon). Tensoring to the abelian
cover $\Lam_w:=\Z[H_1(G_w)]$ (for coprime type $\Lam_w=\Z[t^{\pm1}]$) yields the \emph{Alexander
chain complex}
\[
   \Lam_w\xrightarrow{\,(S,\,T)\,}\Lam_w^2
         \xrightarrow{\,(t^{-(E-o)}-1,\ t^{o}-1)\,}\Lam_w,\qquad
   \partial_1\partial_2=0\ \text{(Lemma~\ref{lem:cf})},
\]
with $H_0=\Z$, $H_2=0$, and $H_1=A_w$ the Alexander module, of order $R_w$ --- the homological
mirror of $H^2=M/(S,T)$ in Theorem~\ref{thm:alex} (Crowell). Since $R_w\ne0$ the complex is
acyclic over $\Q(t)$, hence carries a Reidemeister--Milnor torsion; by Milnor's formula
\[
   \tau(G_w,\chi)\;\doteq\;\frac{R_w(t)}{t-1}\ \in\ \Q(t)^\times/\!\pm t^{\Z}
\]
(the $t-1$ being the order of $H_0=\Lam_w/(t-1)$). Thus $R_w$ is the order of $H_1$ and the
numerator of the torsion --- the master invariant of which the catalogue computes the polynomial
part. (Here $H_2=\ker(S,T)=0$, as $\Lam_w$ is a domain and $S,T\ne0$; and the ascending chain
of Fitting ideals of $A_w$, whose first term is the order ideal $(R_w)$ rationally, gives
strictly finer invariants than the polynomial alone.)
\end{rem}

\begin{rem}[The non-coprime stratum is torsion in $H_1$]\label{rem:noncoprime}
The first homology $H_1(G_w)=G_w^{\mathrm{ab}}=\Z^2/\langle(o,E-o)\rangle$ is read off the Smith
normal form of the relator vector $[\,o,\ E-o\,]$, namely $[\,d,\ 0\,]$ with $d=\gcd(E,o)$:
\[
   H_1(G_w)\;\cong\;\Z\oplus\Z/d .
\]
Coprime type is exactly $d=1$ --- the relator class $(o,E-o)$ primitive, $H_1$ free of rank one
on $t$. For $d>1$ the $\Z/d$ summand is the formal variable $s$ of \S1--\S2 with $s^{d}=1$, and
\[
   \Q\otimes\Z[H_1(G_w)]\;=\;\Q[t^{\pm1},s]/(s^{d}-1)\;\cong\;\prod_{e\mid d}\Q(\zeta_e)[t^{\pm1}],
\]
which is exactly the Chinese-remainder splitting of \S2. So the $s$-grading is the torsion of the
relator's homology class --- not a second free variable, since $H_1$ has rank one. Accordingly
$R_w$ is the $e=1$ (trivial-character) component of the $\Z/d$-twisted Alexander module over
$\Z[H_1(G_w)]=\Z[t^{\pm1},s]/(s^{d}-1)$, and specialising $s\mapsto\zeta_e$ --- the very programme
of \S1 --- is its isotypic decomposition over the $\Z/d$-cover, the components $e>1$ being the
genuinely new data of the non-coprime stratum.
\end{rem}

\begin{rem}[Reciprocity and the duality locus]\label{rem:duality}
For a knot the Alexander polynomial is symmetric, $\Delta(t)\doteq\Delta(t^{-1})$, forced by
Poincar\'e duality of the complement and refined by a Blanchfield linking pairing
$A_w\times A_w\to\Q(t)/\Lam_w$. \emph{This fails here in general}: the cores are usually not
self-reciprocal --- e.g.\ the $-17$ core $t^{11}-t^{10}+t^{7}-t^{6}+t^{4}-t+1$ is not --- because
the $2$-complex of $G_w$ is not a closed manifold. What survives is the \emph{external} symmetry
of Proposition~\ref{prop:cores}: complement-then-reverse sends $w$ to a word with the same core,
$R_{\overline{w^{R}}}=R_w$, via $S\leftrightarrow T$ and $t\mapsto t^{-1}$. The self-reciprocal
cores --- those fixed by reciprocation, $R_w(t)\doteq R_w(t^{-1})$ --- form a distinguished
\emph{duality locus}: of the $1529$ distinct cores with $E\le16$, exactly $107$ are
self-reciprocal --- the $55$ cyclotomic cores (all reciprocal) and $52$ non-cyclotomic ones. The
latter are the reciprocal cores, every \emph{Salem} core among them (Salem polynomials are
reciprocal), the smallest being Lehmer's polynomial; the non-reciprocal cores, like the $-17$
core, lie off the locus. A sufficient condition to land on it is that $w$ be a palindrome up to
rotation ($w^{R}\sim w$): then $R_w=R_{w^{R}}\doteq R_w(t^{-1})$. Whether the words on this locus
carry an honest Blanchfield-type sesquilinear pairing on $A_w$ --- a true duality --- is the
natural open question.
\end{rem}


\end{document}
